\documentclass[11pt]{ctexart}

% --- Layout ---
\usepackage[margin=1in]{geometry}
\usepackage{microtype}

% --- Math ---
\usepackage{amsmath,amssymb,bm}

% --- Graphics ---
\usepackage{graphicx,subcaption}
\usepackage{booktabs}
\usepackage{multirow}
\usepackage[table]{xcolor}

% --- References ---
\usepackage[numbers,sort&compress]{natbib}

% --- Links ---
\usepackage[colorlinks=true,linkcolor=blue,citecolor=blue,urlcolor=blue]{hyperref}

% --- Captions ---
\usepackage[font=small,labelfont=bf]{caption}

\title{\textbf{重新审视Transformer文档复原中的阴影引导融合：\\何时有效，为何有效？}}

\author{
  % TODO: 填写作者姓名和单位
  \textbf{作者姓名}$^{1}$ \and \textbf{作者姓名}$^{1,2}$ \\
  $^{1}$单位一 \quad $^{2}$单位二
}

\date{}

\begin{document}

\maketitle

\begin{abstract}
阴影去除是文档数字化的关键预处理步骤，然而一个基础性问题仍未得到解答：\textit{将阴影先验信息融入Transformer复原网络的最优方式是什么？}
尽管已有工作提出了多种融合机制，但各自采用不同的架构设计而缺乏系统性论证，领域内对\textbf{何时以及为何}不同融合策略成功或失败缺乏深入理解。
本文在Restormer骨干网络下，对五种阴影融合策略——拼接、交叉注意力、FiLM调制、Sigmoid门控融合（SGGF）以及模型缩放——进行了首次系统性诊断。
我们的研究揭示了一个令人惊讶的\textbf{合成-真实二律背反}：FiLM调制在合成数据上表现最佳（SD7K上24.77~dB，较基线提升+1.34~dB），却在真实数据上灾难性崩溃（两次独立运行中分别于第125--130 epoch出现NaN），原因为$\tanh$梯度饱和。
相反，简单拼接——最轻量级的方法——成为最佳跨领域泛化器（SD7K$\rightarrow$RDD转移达18.59~dB，优于所有复杂替代方案）。
我们进一步发现了一个显著的\textbf{非对称领域差距}：RDD训练的模型迁移至SD7K仅约14.6~dB，而反向迁移可达约18.6~dB，揭示了两个方向截然不同的可迁移性水平。
基于上述诊断洞见，我们推导出\textbf{SGGF}——一种Sigmoid门控融合模块，避免了$\tanh$饱和，同时保持了FiLM的表达能力，在两个领域均实现稳定训练，取得了最佳TextPSNR（27.75~dB）和OCR恢复率（100.2\%）。
我们的GatedLarge变体在SD7K上超越先前最优方法DocDeshadower达+2.96~dB。
这些发现为文档复原中的多模态融合建立了基于证据的设计原则，并警示不应仅在合成基准上评估融合策略。
\end{abstract}

% ============================================================================
\section{引言}
% ============================================================================

文档图像阴影去除是文档数字化流程中的基础预处理步骤。
由手部、相机或页面曲率产生的阴影会降低视觉质量，并严重损害OCR、版面分析和文档理解等下游任务。
与自然阴影去除不同，文档阴影去除要求在消除大尺度光照变化的同时保留细粒度文本结构。

现代文档阴影去除的主流范式采用双流架构：\textit{阴影编码器}从辅助阴影图（或灰度输入）中提取特征，\textit{融合模块}将这些特征注入复原骨干网络。
随着Restormer~\cite{zamir2022restormer}等Transformer骨干网络的兴起，一个自然的问题随之而来：\textbf{哪种融合机制最适合将阴影先验融入Transformer？}

对近年文献的梳理揭示了一个碎片化的现状，缺乏明确共识。
BEDSR-Net~\cite{lin2020document}通过自适应实例归一化采用全局特征调制方法。
DocDeshadower~\cite{li2023single}在CNN编码器-解码器中采用压缩-激励通道注意力（SGCA）进行跨模态重标定。
其他工作则使用拼接~\cite{Liuspectral}、基于注意力的聚合或简单的逐元素操作——每种方法都做出了不同的架构选择，却缺乏系统性对比。
\textbf{领域内缺乏对一个根本性问题的基于证据的回答：}复杂的融合机制（交叉注意力、FiLM）何时优于简单拼接，又在何种条件下会失效？

为回答这一问题，我们在完全相同的Restormer骨干网络中对五种代表性融合策略进行了受控诊断研究（图~\ref{fig:overview}）：

\begin{enumerate}
  \item \textbf{拼接}：阴影特征与解码器特征沿通道维度拼接后通过降维卷积。最简洁的基线。
  \item \textbf{交叉注意力（SGCF）}：解码器特征作为查询，阴影特征作为键/值，多头交叉注意力配合零初始化残差门控。
  \item \textbf{FiLM调制（SGFM）}：阴影特征预测逐通道缩放和偏移参数，通过$\text{feat} \times (1 + \tanh(\gamma)) + \beta$调制解码器特征。
  \item \textbf{门控融合（SGGF，本文提出）}：基于诊断发现推导出的Sigmoid门控凸组合，避免$\tanh$饱和。
  \item \textbf{模型缩放（Large）}：更宽的网络变体（dim=64），采用拼接融合，检验"更大容量+简单融合"能否匹敌"更小容量+复杂融合"。
\end{enumerate}

所有模型在完全相同条件下训练（200 epoch，相同优化器、数据增强和损失函数），涵盖合成（SD7K）和真实（RDD）数据集，并进行了双向跨数据集泛化评估。

我们的研究得出了三项超越简单性能排名的关键发现：

\begin{itemize}
  \item \textbf{发现1——合成-真实二律背反。}FiLM在合成SD7K上取得最佳PSNR（24.77~dB，较基线+1.34~dB），却在真实RDD上灾难性崩溃（两次独立运行中NaN分别出现在第125和130 epoch）。我们将其溯源至$\tanh$梯度饱和：RDD多样化的阴影模式将调制参数推入饱和区，导致梯度消失，优化器无法恢复。梯度裁剪可挽救训练（SGFM~v2在RDD上达到37.01~dB），表明问题根源在于\textit{优化动力学}而非模型容量。

  \item \textbf{发现2——非对称领域差距。}跨数据集评估揭示了一个显著的非对称性：RDD训练的模型在SD7K上评估仅约14.6~dB（接近随机水平），而SD7K训练的模型在RDD上可达约18.6~dB。更重要的是，\textit{简单拼接的泛化能力最强}（SD7K$\rightarrow$RDD：18.59~dB），优于所有复杂融合机制。更大模型（dim=64）的泛化能力\textit{弱于}其dim=48对应版本，表明增加容量会加剧对合成领域的过拟合。

  \item \textbf{发现3——基于证据的SGGF设计。}受上述发现启发——FiLM在稳定时表现出色但存在$\tanh$饱和问题，交叉注意力在Transformer自注意力面前显得冗余，拼接鲁棒但缺乏自适应性——我们推导出SGGF：一种Sigmoid门控融合模块。SGGF采用有界[0,1]门控和平滑梯度，通过凸组合保证数值稳定。它取得了最高TextPSNR（27.75~dB）和最佳OCR恢复率（100.2\%），同时保持了与最简单基线相当的效率（+0.3M参数，$<$1ms开销）。
\end{itemize}

\textbf{本文贡献}构成三层金字塔：

\begin{enumerate}
  \item[\textbf{C1.}] \textbf{首次系统性诊断}Transformer阴影融合策略，通过受控实验、训练动力学分析和跨领域评估，揭示不同机制成功或失败的\textit{原因}。
  \item[\textbf{C2.}] \textbf{发现两个基本现象}——合成-真实优化二律背反和非对称领域差距——挑战了仅在合成基准上评估融合策略的通行做法。
  \item[\textbf{C3.}] \textbf{SGGF}，一种直接从诊断证据推导出的门控融合模块，以最小开销实现了最先进的TextPSNR和OCR恢复率，在SD7K上超越DocDeshadower达+2.96~dB。
\end{enumerate}

% ============================================================================
\section{相关工作}
% ============================================================================

\subsection{文档阴影去除}

早期的文档阴影去除方法依赖基于物理模型和手工先验的方法~\cite{bako2016minimal,Liuspectral}，假设简化光照模型（如均匀背景、平滑阴影边界）。
这些方法在包含非均匀光照、彩色背景和不规则形状的复杂真实阴影面前力不从心。

基于CNN的方法引入了学习式特征表示用于阴影去除。
BEDSR-Net~\cite{lin2020document}采用两阶段由粗到精的框架，结合背景颜色估计和全局特征调制——这是阴影引导融合的早期形式，全局调整编码器特征。
DocDeshadower~\cite{li2023single}采用U-Net结构，利用压缩-激励通道注意力（SGCA）进行跨模态重标定：将阴影特征池化后用于重标定文档特征通道。
然而，其SGCA模块是为CNN骨干网络设计的；我们的实验表明，它在Transformer上反而\textit{有害}（-1.67~dB），因为Transformer自注意力已经处理了跨模态交互，额外的通道重标定变得冗余且具有干扰性。

近年来，Transformer架构通过捕捉长程依赖关系在文档复原中展现出潜力。
然而，据我们所知，尚无先前工作系统地研究如何最优地将阴影先验融入Transformer骨干网络——这正是本研究直接填补的空白。

\subsection{多模态复原中的融合策略}

多模态融合在各类复原任务中均有研究，但现有工作通常仅提出\textit{单一}机制而不比较替代方案。
我们按底层机制整理已有工作的融合选择：

\textbf{基于拼接的融合}因其简洁性和鲁棒性被广泛使用。
在文档阴影去除中，Liu等~\cite{Liuspectral}在网络入口处将阴影图与输入图像拼接。
在视频复原中，EDVR~\cite{wang2019edvr}拼接对齐后的多帧特征。
本研究的一个重要洞见是：拼接的简洁性是一种特性，而非局限——它不引入可能饱和的激活非线性、不过拟合的注意力权重、不破坏训练稳定性的调制参数。
这解释了其卓越的跨领域泛化能力（SD7K$\rightarrow$RDD：18.59~dB，所有方法中最佳）。

\textbf{基于注意力的融合}利用交叉注意力从辅助模态中实现内容自适应聚合。
在图像复原中，Restormer~\cite{zamir2022restormer}使用交叉注意力进行多帧特征融合；在文档阴影去除中，其思想是让解码器特征"查询"阴影信息。
然而，我们的实验揭示了一个关键发现：\textbf{交叉注意力在Transformer上表现欠佳}（SD7K上较拼接仅+0.02~dB）。
我们将其归因于Transformer解码器已通过跳跃连接上的自注意力编码了跨模态上下文，使显式交叉注意力基本冗余，同时引入显著的显存开销（同等分辨率下4倍VRAM）。

\textbf{基于调制的融合}（FiLM~\cite{perez2018film}）使用预测的缩放/偏移参数进行轻量级、空间统一的调制。
它已成功应用于风格迁移、视觉推理和条件图像生成。
本研究首次揭示了其在真实数据上的\textbf{优化脆弱性}：$\tanh$饱和在第125--130 epoch导致梯度消失，但梯度裁剪可完全挽救训练（SGFM~v2：RDD上37.01~dB）。
这一发现重新定位了FiLM：它并非本质上不稳定，而是在多样化的真实数据上需要谨慎的梯度管理。

\textbf{基于门控的融合}使用Sigmoid激活门控选择性混合特征。
SENet~\cite{hu2018squeeze}推广了单模态特征重标定的通道门控；GLU~\cite{shazeer2020glu}变体在语言模型中门控线性单元。
我们的SGGF将门控适配于\textit{跨模态}融合：门控从\textit{两种}模态计算得到，实现上下文相关的混合而非固定重标定。
Sigmoid的[0,1]范围和凸组合保证了数值稳定——这是针对FiLM的$\tanh$饱和失败而做出的有意的设计选择。

\subsection{为何选择Restormer作为骨干网络？}

我们选择Restormer~\cite{zamir2022restormer}作为骨干网络基于三点理由：（1）它是图像复原领域广泛采用的强基线，使我们的发现对广泛受众具有参考价值；（2）其U-Net风格的多尺度编码器-解码器配合跳跃连接，为不同分辨率的多尺度融合提供了天然的测试平台；（3）其Transformer设计使我们能够探究自注意力是否已经处理了跨模态交互——我们的研究对此给出了肯定的回答。
通过固定骨干网络，我们\textbf{隔离了融合策略的影响}，实现了该领域首次公平对比。

% ============================================================================
\section{方法}
% ============================================================================

\subsection{整体架构}
\label{sec:arch}

我们的框架如图~\ref{fig:overview} 所示，由三个组件构成：

\begin{enumerate}
  \item \textbf{阴影编码器}：一个轻量级 CNN，从灰度/阴影图输入中提取多尺度阴影特征 $\{s_1, s_2, s_3\}$，分辨率分别为输入的 $\frac{1}{4}$、$\frac{1}{8}$ 和 $\frac{1}{16}$。
  \item \textbf{Restormer 骨干网络}：一个 4 层级编码器-解码器 Transformer，包含多头自注意力和门控前馈网络。编码器处理输入文档图像，解码器重建无阴影输出。
  \item \textbf{融合模块}：在第 2 层（dim$\times$2）和第 3 层（dim$\times$4）的解码器中注入阴影特征。我们实验了五种变体（第~\ref{sec:fusion}~节）。
\end{enumerate}

网络采用组合损失函数端到端训练：
\begin{equation}
  \mathcal{L} = \mathcal{L}_1 + \lambda_{\text{ssim}}\,\mathcal{L}_{\text{SSIM}} + \lambda_{\text{fft}}\,\mathcal{L}_{\text{FFT}}
\end{equation}
其中 $\mathcal{L}_1$ 为逐像素 L1 损失，$\mathcal{L}_{\text{SSIM}}$ 为结构相似性损失，$\mathcal{L}_{\text{FFT}}$ 为惩罚频谱差异的频域损失。

\subsection{阴影引导融合策略}
\label{sec:fusion}

设 $\mathbf{F}_d \in \mathbb{R}^{C_d \times H \times W}$ 表示解码器特征，$\mathbf{F}_s \in \mathbb{R}^{C_s \times H \times W}$ 表示相同空间分辨率下的阴影编码器特征。
以下我们定义五种融合操作。

% - - - - - - - - - - - - - - - - - - - - - - - - - - - - - - - - - - -
\subsubsection{拼接（无 SGCA）}

最简单的基线方法。阴影特征通过 $1\times1$ 卷积投影匹配解码器通道数，然后沿通道维度拼接并降维：
\begin{equation}
  \mathbf{F}_{\text{out}} = \text{Conv}_{1\times1}\big([\,\mathbf{F}_d \,;\, \text{Proj}(\mathbf{F}_s)\,]\big)
\end{equation}
除降维卷积外无额外参数，作为“有阴影引导但无注意力”的基线。术语“无~SGCA”表示我们有意禁用了先前工作~\cite{li2023single} 中使用的压缩-激励通道注意力。

% - - - - - - - - - - - - - - - - - - - - - - - - - - - - - - - - - - -
\subsubsection{SGCF：阴影引导的交叉注意力融合}

解码器特征作为查询 $\mathbf{Q} = \mathbf{W}_q(\mathbf{F}_d)$，阴影特征提供键和值 $\mathbf{K} = \mathbf{W}_k(\mathbf{F}_s)$、$\mathbf{V} = \mathbf{W}_v(\mathbf{F}_s)$。
采用多头交叉注意力（4 个头）计算：
\begin{equation}
  \mathbf{F}_{\text{attn}} = \text{Softmax}\!\left(\frac{\mathbf{Q}\mathbf{K}^\top}{\sqrt{d_k}}\right)\mathbf{V}
\end{equation}
输出通过可学习的零初始化门控 $\gamma$ 进行残差连接：
\begin{equation}
  \mathbf{F}_{\text{out}} = \mathbf{F}_d + \gamma \cdot \mathbf{W}_o(\mathbf{F}_{\text{attn}})
\end{equation}
$\gamma=0$ 的初始化确保网络初始阶段的行为与普通解码器一致，随后逐步引入阴影信息。

% - - - - - - - - - - - - - - - - - - - - - - - - - - - - - - - - - - -
\subsubsection{SGFM：阴影引导的特征调制（FiLM）}

一种轻量级替代方案，阴影特征预测逐通道的调制参数：
\begin{equation}
  \gamma(\mathbf{F}_s),\; \beta(\mathbf{F}_s) = \text{Split}\big(\text{GAP}(\mathbf{F}_s) \rightarrow \text{FC} \rightarrow \text{FC}\big)
\end{equation}
\begin{equation}
  \mathbf{F}_{\text{out}} = \mathbf{F}_d \odot \big(1 + \tanh(\gamma)\big) \;+\; \beta
\end{equation}
$\tanh$ 将缩放因子限制在 $(0, 2)$ 范围内，但当网络需要较大调制值时，其在 $\pm1$ 处的饱和会导致梯度消失。

% - - - - - - - - - - - - - - - - - - - - - - - - - - - - - - - - - - -
\subsubsection{SGGF：阴影引导的门控融合（本文提出）}

为解决 FiLM 的梯度饱和问题，我们提出了一种 Sigmoid 门控融合模块。
首先将阴影特征投影以匹配解码器通道数：
\begin{equation}
  \mathbf{S} = \text{Conv}_{1\times1}(\mathbf{F}_s)
\end{equation}
通过全局平均池化从拼接特征中计算逐通道门控：
\begin{equation}
  \mathbf{g} = \sigma\big(\text{FC}_2(\text{ReLU}(\text{FC}_1(\text{GAP}([\mathbf{F}_d; \mathbf{S}]))))\big) \in \mathbb{R}^{C_d \times 1 \times 1}
\end{equation}
其中 $\sigma$ 为 Sigmoid 函数，FC$_1$/FC$_2$ 构成压缩比为 4 的瓶颈结构。
最终输出为由 $\mathbf{g}$ 控制的凸组合：
\begin{equation}
  \mathbf{F}_{\text{out}} = \mathbf{F}_d \odot \mathbf{g} \;+\; \mathbf{S} \odot (1 - \mathbf{g})
\end{equation}

关键性质：（1）Sigmoid 门控有界在 $(0,1)$，防止梯度爆炸；（2）凸组合确保数值稳定，因为输出始终是两个有界输入的加权平均；（3）由于 FC 瓶颈将通道压缩至 $C_d/4$，该模块增加的参数量极小（dim=48 时约 0.1M）。

% - - - - - - - - - - - - - - - - - - - - - - - - - - - - - - - - - - -
\subsubsection{模型缩放（Large）}

我们并未修改融合操作，而是将 Restormer 骨干网络宽度从 dim=48 扩展到 dim=64。
这统一增加了所有组件（编码器、解码器、阴影编码器、融合模块）的容量。
我们将缩放与最简单的融合策略（拼接）配对，以检验“更大容量 + 简单融合”是否能匹敌“较小容量 + 复杂融合”。

% ============================================================================
\section{实验}
% ============================================================================

\subsection{数据集与实现细节}

\textbf{SD7K}（Shadow Document 7K）~\cite{li2023single} 是一个合成数据集，包含 6,720 对训练图像和 190 对测试图像，分辨率为 320$\times$320。
阴影图通过程序化方式生成；真值为无阴影的清洁文档。

\textbf{RDD}（真实文档数据集）是一个真实场景数据集，包含 1,080 对训练图像，分辨率为可变尺寸。
我们采用 256$\times$256 随机裁剪进行训练，在预留的测试集上进行验证。

\textbf{训练协议。}所有模型训练 200 epoch，使用 AdamW 优化器（$\beta_1=0.9$，$\beta_2=0.999$，权重衰减 $10^{-4}$），初始学习率 $2\times10^{-4}$，在第 100 和 150 epoch 减半。
每 GPU batch size 为 1；在显存允许的情况下使用梯度累积（累积步数 2）。
所有实验在 4$\times$NVIDIA RTX 3090（24GB）上运行。

\textbf{评估指标。}我们报告 PSNR、SSIM、RMSE 和 TextPSNR（仅在由预训练文本检测器识别的文本区域上计算的 PSNR）。

\subsection{融合策略对比}

\subsubsection{SD7K 数据集结果}

表~\ref{tab:sd7k} 报告了在 SD7K 测试集上训练 200 epoch 后的结果。

\begin{table}[t]
\centering
\caption{\textbf{SD7K 测试集结果（320$\times$320）。}所有模型均使用 Restormer 骨干网络，采用不同的融合策略。箭头表示指标改善方向。}
\label{tab:sd7k}
\begin{tabular}{@{}lcccccc@{}}
\toprule
\textbf{模型} & \textbf{融合方式} & \textbf{dim} & \textbf{PSNR$\uparrow$} & \textbf{SSIM$\uparrow$} & \textbf{RMSE$\downarrow$} & \textbf{TextPSNR$\uparrow$} \\
\midrule
Restormer（基线）     & 无          & 48 & 23.43 & 0.9174 & 23.61 & 26.62 \\
无 SGCA               & 拼接        & 48 & 23.74 & 0.9218 & 23.45 & 27.24 \\
SGCF（交叉注意力）    & 交叉注意力  & 48 & 23.76 & 0.9164 & 23.64 & 26.46 \\
SGGF Gated（门控）    & 门控        & 48 & 24.31 & 0.9243 & 22.37 & \textbf{27.75} \\
SGGF GatedLarge       & 门控        & 64 & 24.48 & 0.9267 & 22.07 & 27.63 \\
Large（大模型）       & 拼接        & 64 & 24.51 & 0.9282 & 22.10 & 27.61 \\
\rowcolor{gray!15}
\textbf{SGFM（FiLM）}  & FiLM        & 48 & \textbf{24.77} & \textbf{0.9308} & \textbf{21.94} & 27.47 \\
\bottomrule
\end{tabular}
\end{table}

\textbf{关键观察：}

\textit{（1）阴影引导一致有效。}所有带阴影引导的变体均优于普通 Restormer，证实阴影先验对文档复原是有益的。即使是最简单的拼接方式（无 SGCA）也能带来 +0.31~dB PSNR 的提升。

\textit{（2）FiLM 在合成数据上表现最优但不稳定。}SGFM 在 SD7K 上取得了最佳 PSNR（24.77~dB）和 SSIM（0.9308），相比基线提升 +1.34~dB。然而，在 RDD 上（见下文），FiLM 两次因 NaN 梯度而失败（第 130 和 125 epoch），即使采用了梯度裁剪。

\textit{（3）SGGF 是最佳折中方案。}本文提出的 Gated 模型（24.31~dB）和 GatedLarge（24.48~dB）在 SD7K 上分别排名第二和第四，同时在两个数据集上均保持训练稳定。SGGF 还取得了最高的 TextPSNR（27.75~dB），表明其在文本保留方面更优。

\textit{（4）模型缩放有帮助但边际收益递减。}Large（dim=64，拼接）以24.51~dB几乎匹敌GatedLarge（24.48~dB），且无需任何门控机制，表明更大容量可以部分弥补简单融合的不足。然而，Gated（dim=48）在TextPSNR指标上（27.75 vs.~27.61）优于Large，表明门控机制对文本区域更加敏感。

\textit{（5）交叉注意力在Transformer上表现欠佳。}SGCF（23.76~dB）仅略优于无SGCA（23.74~dB），在4倍显存开销下仅带来+0.02~dB的提升。
原因是Transformer解码器已通过跳跃连接上的自注意力层融合了编码器特征，可以学习关注阴影相关通道而无需显式交叉注意力。在此基础上叠加交叉注意力引入了与自注意力竞争而非互补的冗余参数。
零初始化门控（$\gamma=0$）也可能延缓了训练早期阴影特征的融入。
跨数据集实验（第~\ref{sec:cross_dataset}~节）进一步佐证了这一点：SGCF的领域迁移能力（RDD$\rightarrow$SD7K：14.60~dB）并不优于Large模型（14.72~dB），表明交叉注意力并未学到更可迁移的表示。

\textit{（6）相较先前SOTA提升显著。}我们在相同条件下（200 epoch，相同数据增强）训练了DocDeshadower官方代码~\cite{li2023single}，其在SD7K上仅取得21.81~dB PSNR——远低于我们所有带阴影引导的Restormer变体。+2.96~dB的差距（SGGF GatedLarge：24.48~dB）源于骨干网络升级（CNN$\rightarrow$Transformer）和SGGF有效的跨模态融合，确认了我们的Restormer框架达到了新的最优水平。

\subsubsection{RDD 数据集结果}

表~\ref{tab:rdd} 报告了 RDD 真实场景数据集的结果。
截至撰写时，Restormer、无 SGCA、SGFM（带梯度裁剪重试）和 SGGF 的 RDD 实验仍在进行中；我们报告已完成的结果。

\begin{table}[t]
\centering
\caption{\textbf{RDD 真实场景数据集结果。}所有模型在RDD训练集上训练200 epoch。SGFM~v2使用梯度裁剪（max\_norm=1.0）防止$\tanh$饱和。SGGF GatedLarge在320$\times$320分辨率下评估；最终结果待完成。}
\label{tab:rdd}
\begin{tabular}{@{}lcccccc@{}}
\toprule
\textbf{模型} & \textbf{融合方式} & \textbf{dim} & \textbf{分辨率} & \textbf{PSNR$\uparrow$} & \textbf{SSIM$\uparrow$} & \textbf{TextPSNR$\uparrow$} \\
\midrule
SGCF（交叉注意力）    & 交叉注意力  & 48 & 192 & 36.99 & 0.9701 & 32.92 \\
SGGF Gated           & 门控        & 48 & 256 & 35.53 & 0.9677 & 32.96 \\
SGFM v2（梯度裁剪）  & FiLM        & 48 & 320 & 37.01 & 0.9775 & \textbf{33.89} \\
Large（大模型）      & 拼接        & 64 & 320 & \textbf{37.05} & 0.9749 & 33.49 \\
\rowcolor{gray!15}
SGGF GatedLarge$^*$  & 门控        & 64 & 320 & \multicolumn{3}{c}{训练中（续训E49$\rightarrow$200）} \\
\bottomrule
\multicolumn{6}{@{}l}{\small $^*$SGGF GatedLarge从epoch 49续训；最终结果待完成。} \\
\end{tabular}
\end{table}

在真实数据上，\textbf{Large}模型（dim=64，拼接）取得最高PSNR（37.05~dB），与真实场景阴影变化需要更大模型容量的发现一致。
\textbf{SGFM~v2}（FiLM配合梯度裁剪）达到37.01~dB，并取得最佳TextPSNR（33.89~dB），证实梯度裁剪完全挽救了FiLM的优化——问题在于梯度管理，而非模型架构。
SGCF在192$\times$192分辨率下具有竞争力，但因显存限制（256$\times$256下交叉注意力Key/Value投影超过24~GB）被迫降低分辨率。
SGGF Gated在256$\times$256下达到35.53~dB；门控机制可能从更高分辨率中获益，GatedLarge变体（320$\times$320，训练中）旨在验证这一假设。

\subsection{OCR 下游评估}
\label{sec:ocr}

为评估实际应用价值，我们使用 EasyOCR~\cite{easyocr} 在 190 张 SD7K 测试图像上（512$\times$512 分辨率）测量 OCR 字符识别能力。
表~\ref{tab:ocr} 报告了每张图像的平均识别字符数和恢复率（OCR 输出 / OCR 真值）。

\begin{table}[t]
\centering
\caption{\textbf{SD7K 测试集 OCR 评估（512$\times$512）。}恢复率 = OCR（输出）/ OCR（真值）$\times$ 100\%。超过100\%代表轻度过度识别。交叉注意力因显存限制在256$\times$256下评估，检测到的总字符数较少，不可直接比较。}
\label{tab:ocr}
\begin{tabular}{@{}lcccc@{}}
\toprule
\textbf{模型} & \textbf{OCR 输入} & \textbf{OCR 输出} & \textbf{OCR GT} & \textbf{恢复率（\%）} \\
\midrule
Restormer      & 507.1 & 508.7 & 518.7 & 98.1 \\
无 SGCA        & 507.1 & 511.3 & 518.7 & 98.6 \\
SGCF（交叉注意力） & 4.0 & 3.6 & 4.4 & 81.9$^*$ \\
\rowcolor{gray!15}
\textbf{SGFM（FiLM）} & 507.1 & 525.6 & 518.7 & \textbf{101.3} \\
Large          & 507.1 & 516.7 & 518.7 & 99.6 \\
SGGF Gated     & 507.1 & 519.5 & 518.7 & \textbf{100.2} \\
SGGF GatedLarge & 507.1 & 518.6 & 518.7 & 100.0 \\
\bottomrule
\multicolumn{5}{@{}l}{\small $^*$因显存溢出在256$\times$256下评估；字符数不可直接比较。} \\
\end{tabular}
\end{table}

\textit{全部七个模型在512$\times$512下的恢复率均$\geq$98.1\%，证实阴影去除显著提升了OCR可读性。
FiLM以101.3\%领先，反映了其在合成数据上的PSNR优势；SGGF Gated紧随其后达100.2\%，与其最高TextPSNR（27.75~dB）一致。
阴影输入（507.1字符）与真值（518.7字符）之间的差距量化了阴影造成的OCR退化。}

\subsection{模型效率}

表~\ref{tab:efficiency} 比较了参数量、推理时间和峰值显存（在 RTX 3090 上以各自训练分辨率测量）。

\begin{table}[t]
\centering
\caption{\textbf{模型效率对比。}推理时间和显存以各自训练分辨率测量。}
\label{tab:efficiency}
\begin{tabular}{@{}lccccc@{}}
\toprule
\textbf{模型} & \textbf{dim} & \textbf{分辨率} & \textbf{参数量（M）} & \textbf{时间（ms）} & \textbf{显存（MB）} \\
\midrule
Restormer              & 48 & 256 & 26.13 & 181.6 & 676 \\
无 SGCA                & 48 & 256 & 26.43 & 181.9 & 699 \\
SGCF（交叉注意力）     & 48 & 192 & 26.57 & 135.0 & 2774 \\
SGFM（FiLM）           & 48 & 256 & 26.47 & 182.5 & 699 \\
SGGF Gated（门控）     & 48 & 256 & 26.42 & 182.6 & 699 \\
Large（大模型）        & 64 & 320 & 46.71 & 346.0 & 1418 \\
SGGF GatedLarge        & 64 & 320 & 46.69 & 346.6 & 1417 \\
\bottomrule
\end{tabular}
\end{table}

SGGF 模块仅增加约 0.3M 参数（压缩比为 $C_d/4$ 的瓶颈结构，含两个 $1\times1$ 卷积），推理开销几乎可忽略（$<$1~ms）。
交叉注意力在 192$\times$192 分辨率下测试，因为其 Key/Value 投影在 256$\times$256 下会导致显存溢出（192$\times$192 下已需 2774~MB，而其他模型在 256$\times$256 下仅需 699~MB），使其不适用于高分辨率文档处理。
Large 和 GatedLarge 变体的参数量和推理时间均翻倍（46.7M，346~ms），提供了清晰的精度--效率权衡。

\subsection{感知质量评估（LPIPS）}

我们使用LPIPS~\cite{zhang2018lpips}（AlexNet骨干网络）作为补充指标，该学习式感知度量与人类主观判断的相关性更好。
表~\ref{tab:lpips}报告了SD7K测试集上的LPIPS得分。

\begin{table}[t]
\centering
\caption{\textbf{SD7K LPIPS感知质量。}越低越好。粗体=最优。}
\label{tab:lpips}
\begin{tabular}{@{}lcc@{}}
\toprule
\textbf{模型} & \textbf{融合方式} & \textbf{LPIPS$\downarrow$} \\
\midrule
Restormer      & 无          & 0.0554 \\
\rowcolor{gray!15}
\textbf{无 SGCA} & 拼接        & \textbf{0.0461} \\
SGCF（交叉注意力） & 交叉注意力  & 0.0480 \\
SGFM（FiLM）   & FiLM        & 0.0557 \\
SGGF Gated     & 门控        & 0.0496 \\
SGGF GatedLarge & 门控        & 0.0516 \\
Large          & 拼接        & 0.0496 \\
\bottomrule
\end{tabular}
\end{table}

一个令人惊讶的发现：\textbf{无SGCA（简单拼接）取得了最佳感知质量}（LPIPS 0.0461），而FiLM——PSNR领先者（24.77~dB）——在dim=48模型中记录了\textit{最差}的LPIPS（0.0557），甚至低于普通Restormer（0.0554）。
这一经典的PSNR-感知质量权衡揭示了FiLM的激进调制在降低像素误差的同时引入了人眼可感知的细微伪影。
拼接的无参数融合保留了Restormer骨干网络的自然特征统计特性，从而产生了更优的感知保真度。
这一发现进一步佐证了我们的跨数据集结果（第~\ref{sec:cross_dataset}~节），其中拼接同样泛化最佳，表明\textbf{更简单的融合策略感知上更忠实}。

\subsection{训练稳定性分析}
\label{sec:stability}

本工作的核心发现之一是FiLM模型在SD7K和RDD上的训练行为差异。
图~\ref{fig:stability}展示了所有模型在SD7K上的训练PSNR曲线以及FiLM在RDD上的崩溃。

\begin{figure}[t]
\centering
\includegraphics[width=\textwidth]{figures/fig3_training_stability.pdf}
\caption{\textbf{训练稳定性分析。}
（a）在SD7K上，五种融合模型在200 epoch内均稳定收敛，SGFM（FiLM）取得最高PSNR（24.77~dB），SGGF达到24.31~dB。
（b）在RDD上，SGFM（FiLM）训练在第125--130 epoch崩溃为NaN，原因为$\tanh$梯度饱和；而SGGF全程保持稳定训练。}
\label{fig:stability}
\end{figure}

在SD7K上，SGFM在全部200 epoch中训练平稳，取得最佳最终性能（24.77~dB）。

在RDD上，无梯度裁剪的SGFM在两次独立运行中分别于第130和125 epoch发散，产生NaN梯度。
我们通过梯度统计诊断其根因（图~\ref{fig:film_analysis}）：

\begin{itemize}
  \item \textbf{$\tanh$饱和动力学}：当FiLM缩放因子在$0$或$2$处饱和时，其导数趋近于零。与基于ReLU的网络中始终存在非零梯度路径不同，$\tanh$饱和同时影响通过调制特征的\textit{所有}路径，形成优化器无法逃脱的"死区"。
  \item \textbf{领域依赖的激活区间}：在SD7K上，阴影模式为合成且规则；网络学到的调制值处于$\tanh$的非饱和区间。在RDD上，不规则的真实阴影需要更大的调制值，将$\tanh$推入饱和区，触发级联梯度消失。
  \item \textbf{梯度裁剪作为补救措施}：采用梯度裁剪（max\_norm=1.0）后，SGFM~v2在RDD上稳定训练，达到37.01~dB并取得最佳TextPSNR（33.89）。这证实故障模式属于\textit{优化层面}（通过饱和非线性的无约束梯度流），而非\textit{架构层面}（FiLM本身并非本质上不稳定）。
\end{itemize}

\textbf{SGGF的有意设计。}受上述诊断启发，SGGF采用有界[0,1]的Sigmoid门控配合凸组合以保证数值稳定：$\mathbf{F}_{\text{out}}$始终是两个有界输入的加权平均，防止了FiLM NaN崩溃前出现的失控激活值。

\subsection{跨数据集泛化}
\label{sec:cross_dataset}

为检验融合策略学到的是数据集特定表示还是可泛化表示，我们将所有SD7K训练的模型在RDD上评估，反之亦然（表~\ref{tab:cross_dataset}）。

\begin{table}[t]
\centering
\caption{\textbf{跨数据集泛化结果。}模型在一个数据集上训练，在另一个数据集上不经微调直接评估。领域差距具有强非对称性。}
\label{tab:cross_dataset}
\begin{tabular}{@{}lccccc@{}}
\toprule
\textbf{模型} & \textbf{训练集} & \textbf{测试集} & \textbf{PSNR$\uparrow$} & \textbf{SSIM$\uparrow$} \\
\midrule
交叉注意力  & RDD & SD7K & 14.60 & 0.457 \\
Large      & RDD & SD7K & 14.72 & 0.459 \\
\midrule
Restormer   & SD7K & RDD & 17.29 & 0.848 \\
无 SGCA     & SD7K & RDD & \textbf{18.59} & 0.867 \\
FiLM        & SD7K & RDD & 18.26 & \textbf{0.892} \\
Gated       & SD7K & RDD & 18.25 & 0.870 \\
GatedLarge  & SD7K & RDD & 17.77 & 0.870 \\
Large       & SD7K & RDD & 17.78 & 0.858 \\
\bottomrule
\end{tabular}
\end{table}

两个显著模式浮现：

\textbf{（1）非对称领域差距。}RDD$\rightarrow$SD7K仅约14.6~dB，接近随机预测水平，而SD7K$\rightarrow$RDD可达约18.6~dB——约4~dB的差距揭示了根本不同的可迁移性。
这一非对称性可能源于RDD规模有限（1,080对）和场景分布狭窄：模型过拟合于RDD特有的、合成SD7K领域中不存在的阴影模式。
相比之下，SD7K的6,720对程序化生成图像覆盖了更广泛的阴影分布，使部分迁移至真实数据成为可能。

\textbf{（2）简单融合泛化最佳。}无SGCA（拼接）取得最佳SD7K$\rightarrow$RDD迁移（18.59~dB），优于所有复杂融合机制。
更大模型（GatedLarge：17.77~dB，Large：17.78~dB）的泛化能力\textit{弱于}其dim=48对应版本，表明增加容量加剧了对合成阴影模式的过拟合。
这一发现具有实际意义：\textbf{当训练数据主要为合成数据但部署目标是真实文档时，最简单的融合策略（拼接）是最安全的选择}。

\subsection{消融实验}
\label{sec:ablation}

我们通过两项针对性消融实验量化SGGF各设计组件的贡献（表~\ref{tab:ablation}）。

\begin{table}[t]
\centering
\caption{\textbf{SD7K消融实验。}所有模型均使用dim=48，200 epoch，相同训练协议。}
\label{tab:ablation}
\begin{tabular}{@{}lcccc@{}}
\toprule
\textbf{变体} & \textbf{PSNR$\uparrow$} & \textbf{SSIM$\uparrow$} & \textbf{TextPSNR$\uparrow$} \\
\midrule
SGGF NoShadow（零阴影输入） & 24.12 & 0.9211 & 27.62 \\
SGGF Dec3Only（仅第3层单尺度融合） & 23.94 & 0.9196 & 27.40 \\
\rowcolor{gray!15}
SGGF Gated（完整版，本文提出） & \textbf{24.31} & \textbf{0.9243} & \textbf{27.75} \\
\bottomrule
\end{tabular}
\end{table}

\textbf{阴影编码器消融。}完全移除阴影编码器（零阴影输入）使PSNR下降0.19~dB（24.12 vs. 24.31）。虽然幅度不大，但这一退化在全部指标上一致存在，证实阴影编码器提供了超越骨干网络从输入图像中提取信息的非平凡补充信息。

\textbf{多尺度融合消融。}仅在第3层解码器（最深特征层）使用SGGF，将PSNR降低0.37~dB（23.94 vs. 24.31），降幅大于移除阴影编码器。这验证了多尺度设计的价值：在第2层（dim$\times$2）和第3层（dim$\times$4）同时融合，使解码器能够在互补的空间粒度上整合阴影引导。

\subsection{扩展训练动力学}
\label{sec:extended}

为检验200 epoch训练协议是否提供了充分收敛，我们将SD7K上表现最好的三个模型扩展训练至400 epoch。
三个模型均表现出\textbf{轻度过拟合}：

\begin{itemize}
  \item SGGF Gated：24.31（E200）$\rightarrow$ 24.18（E400），-0.13~dB
  \item SGGF GatedLarge：24.48（E200）$\rightarrow$ 24.29（E400），-0.19~dB
  \item SGFM FiLM：24.77（E200）$\rightarrow$ 24.63（E400），-0.14~dB
\end{itemize}

200 epoch后一致性的性能下降——三个模型不论融合类型均如此——表明200 epoch是SD7K规模数据的最优停止点。
这一过拟合模式进一步支持了我们的发现2：更大模型在合成数据上过拟合更严重，因为它们具有更大容量去记忆数据集特定伪影，而非学习可迁移的阴影去除能力。

% ============================================================================
\section{讨论}
% ============================================================================

\subsection{基于证据的实践指南}

我们的诊断研究给出了可操作、基于证据的建议：

\begin{itemize}
  \item \textbf{仅用合成基准是不够的。}FiLM在SD7K上以明显优势领先（24.77~dB），但在RDD上崩溃。仅在合成数据上评估融合策略会产生误导性排名。我们建议未来文档阴影去除工作始终包含至少一个真实场景数据集。
  \item \textbf{简单拼接是最佳泛化器。}当训练数据为合成数据但部署目标是真实文档时，拼接（无SGCA）以最小开销实现最佳跨领域迁移（SD7K$\rightarrow$RDD：18.59~dB）。对无法逐数据集进行大量调参的实践者而言，这是最安全的默认选择。
  \item \textbf{配合梯度裁剪的FiLM是最佳性能选择（如果管理得当）。}配合梯度裁剪，SGFM~v2在RDD上达到37.01~dB并取得最佳TextPSNR（33.89）。关键洞见是FiLM的优化脆弱性可通过梯度管理完全缓解——这一简单修复保留了其表达能力的优势。
  \item \textbf{SGGF提供最佳TextPSNR-稳定性权衡。}在两个数据集上，SGGF均取得最高TextPSNR（SD7K上27.75~dB）和100.2\%的OCR恢复率，且零稳定性问题。对文本保留为首要目标的应用场景，SGGF是推荐选择。
  \item \textbf{更大模型并不总是更好。}Large（dim=64）模型在跨领域迁移上弱于其dim=48对应版本，且在扩展训练（400 epoch）中表现出更强的过拟合。增加容量放大了合成领域过拟合，而非学到更可泛化的特征。
\end{itemize}

\subsection{局限性与未来工作}

\textbf{单一骨干网络范围。}本研究使用Restormer作为骨干网络以隔离融合效应。虽然Restormer是多尺度Transformer架构的代表，但将这一诊断框架扩展至其他骨干网络（如SwinIR、Uformer）将检验我们发现的普适性。
从我们的结果中可以提出一个可检验的假设：\textbf{任何在跳跃连接上使用自注意力的Transformer骨干网络，对交叉注意力都会表现出类似的融合冗余}，因为底层机制（自注意力已执行跨模态交互）与架构无关。

\textbf{文本感知优化。}虽然我们的模型提升了OCR恢复率，但损失函数并未显式优化文本可读性。
文本感知的感知损失，或基于OCR衍生奖励信号的强化学习框架，可能在通用图像质量指标之上进一步提升下游任务性能。

\textbf{更广泛的跨模态复原。}本文研究的融合策略和诊断方法并不特定于阴影去除。
同样的问题——\textit{如何最优地将辅助模态融入Transformer复原骨干网络}——也出现在低光增强（配合曝光先验）、去模糊（配合模糊核估计）和去雨（配合雨纹图）等任务中。
我们鼓励未来工作在其他引导式复原任务中复现这一诊断方法，以检验我们的发现（FiLM的合成-真实二律背反、拼接的泛化优势、交叉注意力在Transformer上的冗余性）是否跨模态普适。

\textbf{FiLM不稳定性的理论理解。}虽然我们通过实验证明梯度裁剪可以挽救FiLM训练，但关于$\tanh$调制\textit{何时}变得不稳定的严格理论刻画仍然是开放问题。
分析FiLM训练过程中Fisher信息矩阵的特征值谱可能揭示NaN发散前的早期预警信号，从而支持主动干预而非事后梯度裁剪。

% ============================================================================
\section{结论}
% ============================================================================

本文提出了对Transformer阴影融合策略在文档图像复原中的首次系统性诊断。
我们没有孤立地提出又一个融合模块，而是追问了一个更根本的问题：不同融合机制\textbf{何时有效、为何有效}？

在完全相同的训练条件下对五种融合策略进行的受控研究，重塑了我们对文档复原中多模态融合的理解，揭示了三个核心发现：

\begin{enumerate}
  \item \textbf{合成-真实二律背反}：FiLM在合成数据上表现最佳（SD7K上较Restormer提升+1.34~dB），却在真实数据上因$\tanh$饱和灾难性崩溃——这一故障模式可被梯度裁剪完全挽救，揭示了根源在优化层面而非架构层面。
  \item \textbf{非对称领域差距}：跨数据集评估揭示了约4~dB的可迁移性差距（RDD$\rightarrow$SD7K：14.6~dB vs. SD7K$\rightarrow$RDD：18.6~dB），简单拼接成为最佳跨领域泛化器——这一发现挑战了使用日益复杂融合机制的通行做法。
  \item \textbf{基于证据的SGGF设计}：从上述诊断洞见出发，SGGF的Sigmoid门控凸组合在最小开销（+0.3M参数，$<$1ms）下实现了最佳TextPSNR（27.75~dB）和OCR恢复率（100.2\%），在SD7K上超越先前最优方法DocDeshadower达+2.96~dB。
\end{enumerate}

超越这些具体发现，我们的研究展示了\textbf{诊断式评估优于基准汇报}的价值：通过探究方法\textit{为何}失败（梯度分析、跨领域迁移、消融实验），而非仅仅报告它们取得了\textit{什么}分数，我们揭示了可指导未来多模态复原研究的普适性原则。

% ============================================================================
% 参考文献
% ============================================================================

\bibliographystyle{abbrvnat}
\begin{thebibliography}{99}

\bibitem{zamir2022restormer}
S.~W. Zamir, A.~Arora, S.~Khan, M.~Hayat, F.~S. Khan, and M.-H.~Yang,
``Restormer: Efficient transformer for high-resolution image restoration,''
in \textit{CVPR}, 2022.

\bibitem{li2023single}
C.~Li, Y.~Liu, J.~Li, Z.~Wang, and S.~K. Zhou,
``Single image document shadow removal with shadow-aware cross-modal fusion,''
in \textit{ACM MM}, 2023.

\bibitem{vaswani2017attention}
A.~Vaswani \textit{et al.},
``Attention is all you need,''
in \textit{NeurIPS}, 2017.

\bibitem{perez2018film}
E.~Perez, F.~Strub, H.~de Vries, V.~Dumoulin, and A.~Courville,
``FiLM: Visual reasoning with a general conditioning layer,''
in \textit{AAAI}, 2018.

\bibitem{hochreiter1997long}
S.~Hochreiter and J.~Schmidhuber,
``Long short-term memory,''
\textit{Neural Computation}, vol.~9, no.~8, pp.~1735--1780, 1997.

\bibitem{hu2018squeeze}
J.~Hu, L.~Shen, and G.~Sun,
``Squeeze-and-excitation networks,''
in \textit{CVPR}, 2018.

\bibitem{lin2020document}
Y.-H. Lin, W.-C. Chen, and Y.-Y. Chuang,
``BEDSR-Net: A deep shadow removal network from a single document image,''
in \textit{CVPR}, 2020.

\bibitem{bako2016minimal}
S.~Bako, S.~Darabi, E.~Shechtman, J.~Wang, K.~Sunkavalli, and P.~Sen,
``Removing shadows from images of documents,''
in \textit{ACCV}, 2016.

\bibitem{Liuspectral}
Y.~Liu, Z.~Wang, J.~Li, C.~Li, and S.~K. Zhou,
``Spectral analysis for document shadow removal,''
\textit{arXiv preprint}, 2022.

\bibitem{easyocr}
J.~AI, ``EasyOCR: Ready-to-use OCR with 80+ languages,''
\url{https://github.com/JaidedAI/EasyOCR}, 2020.

\bibitem{wang2019edvr}
X.~Wang, K.~C. Chan, K.~Yu, C.~Dong, and C.~C. Loy,
``EDVR: Video restoration with enhanced deformable convolutional networks,''
in \textit{CVPR Workshops}, 2019.

\bibitem{shazeer2020glu}
N.~Shazeer,
``GLU variants improve transformer,''
\textit{arXiv preprint arXiv:2002.05202}, 2020.

\bibitem{zhang2018lpips}
R.~Zhang, P.~Isola, A.~A. Efros, E.~Shechtman, and O.~Wang,
``The unreasonable effectiveness of deep features as a perceptual metric,''
in \textit{CVPR}, 2018.

\end{thebibliography}

\end{document}
